\section{Mục tiêu nghiên cứu}

Nghiên cứu này được thực hiện nhằm đạt được các mục tiêu cụ thể sau:

Thứ nhất, khảo sát và phân tích thực trạng nhu cầu, những khó khăn về mặt thời gian cùng các công đoạn tiêu tốn nhiều công sức nhất của giảng viên Đại học Bách khoa Hà Nội trong quy trình thiết kế bài giảng và soạn thảo học liệu. Mục tiêu này giúp xác định rõ các điểm nghẽn thực tế mà giảng viên đang đối mặt để định hình giải pháp hỗ trợ phù hợp.

Thứ hai, xác định các nguyên tắc thiết kế và chuẩn sư phạm làm nền tảng khoa học cho việc tích hợp trí tuệ nhân tạo tạo sinh vào dạy học. Nghiên cứu tập trung vào việc áp dụng Thang đo nhận thức Bloom cải tiến trong thiết kế câu hỏi đánh giá và Thuyết kiến tạo trong đề xuất các hoạt động học tập chủ động.

Thứ ba, thiết kế khung chức năng và xây dựng sản phẩm thử nghiệm ứng dụng các mô hình ngôn ngữ lớn để hỗ trợ giảng viên giải quyết các tác vụ soạn bài cụ thể. Công cụ tập trung vào khả năng tự động hóa việc tóm tắt tài liệu chuyên ngành, phân rã cấu trúc bài giảng để chuẩn bị bản trình chiếu và sinh câu hỏi kiểm tra đánh giá theo các mức độ nhận thức.

Thứ tư, đề xuất các khuyến nghị và định hướng ứng dụng thực tế sản phẩm công nghệ nhằm hỗ trợ hiệu quả cho giảng viên trong công tác giảng dạy, đồng thời bảo đảm tính chuẩn hóa sư phạm và tính chính xác chuyên môn của học liệu.
